\documentclass[10pt]{article}
\usepackage[final,main]{neurips_2026}
\usepackage{lmodern}
\usepackage{amsmath,booktabs,geometry,longtable,array,float}
\geometry{margin=0.68in}
\title{Global Context and Feasibility Ablations for PIGNN\\\large GBnetwork and case300}
\author{PIGNN--Attn--LS--PPC experimental report}
\date{28 August 2026}
\begin{document}
\maketitle

\section{Executive summary}

This report consolidates the controlled global-context experiments for
\texttt{case300} and \texttt{GBnetwork}.  The experiments retain the PIGNN
backbone and $K=40$ recurrent corrections and change only the graph-context
sidecar.  G0 is the signed-log PIGNN reference; G1--G3 are global-context
ablations; G4--G5 add dynamic sigmoid gating; and RANGE uses persistent master
tokens.  GraphKit and GridSFM values are included as external reference rows
from the ppNR-v2 four-model report.

The main conclusions are:
\begin{itemize}
\item On \texttt{GBnetwork}, G2 (attention before local processing) gives the
      best voltage result among G1--G5, while G3 (attention after local
      processing) gives the best P/Q residuals among those variants.
\item G4 and G5 do not improve consistently over their ungated counterparts;
      their sigmoid gates remain mostly open and therefore do not provide a
      useful selective safeguard.
\item RANGE is now complete and test-scored.  $M=1$ is the best RANGE setting
      but lands inside the G1--G5 band rather than beyond it, and $M=4$ and
      $M=8$ are clearly worse: their masters collapse to identical
      representations (attention and embedding cosine $+1.000$).  Persistent
      master tokens are therefore not a productive direction for GBnetwork.
\item A per-bus residual breakdown shows PIGNN is \emph{not} broadly worse than
      GridSFM on GBnetwork: G3's median residual is lower on both channels.  The
      entire gap is an extreme tail concentrated on four 22-kV buses
      ($0.18\%$ of the grid), which carry up to $20\%$ of the total $|\Delta Q|$
      mass, and it scales with $|Y_{ii}|$.  GridSFM exceeds 10 pu in no bin of
      any breakdown, which reads as an output bound rather than better typical
      accuracy.
\item Measured against a predictor that ignores the scenario entirely --- each
      bus held at its own training-split mean --- every model here loses on
      $|V|$ RMSE and on both median residuals, and wins only on bus angle.  Two
      thirds of that angle win is gauge drift rather than flow accuracy
      (Section~\ref{sec:metrology}).
\item The negative voltage $R^2_{\sigma}$ throughout this report is a
      \emph{calibration}
      failure, not collapse: after each bus's static profile is removed the
      models still explain 59--75\% of the variance.  Subtracting a per-bus
      offset fitted on the training split --- no retraining --- cuts $|V|$ RMSE
      by $2.0$--$2.7\times$, turns every $R^2_{\sigma}$ positive, and improves median
      $|\Delta Q|$ by $3$--$4\times$.  It does not improve the worst-bus tail.
\item Replacing the log-cosh physics term with MSE (G3-mse) is worse on every
      metric and an order of magnitude worse on the worst bus
      ($\max|\Delta P|$ $18.6\to197.8$ pu), which is the opposite of what a
      quadratic penalty was expected to do to the tail.
\item The GridSFM values 1.10 and 4.81, like the corresponding PIGNN log
      fields, are \emph{means of per-scenario maxima}, not global maxima over
      the complete test set.  They are comparable as reducers when split,
      masks and precision also match.
\end{itemize}

Unless a table says otherwise, $R^2$ in this report means $R^2_{\sigma}$ of
Section~\ref{sec:r2notation}: $1-(\mathrm{RMSE}/\sigma)^2$ against a per-bus
constant predictor.  It is not a tracking measure; the tracking coefficient
$R^2_{\mathrm{bus}}$ appears there.

\section{Experimental protocol}

All G1--G5 and RANGE runs use the same PPC adapter, DC initialization,
train/validation/test split, seed 42 and complex128 parquet decoding.  The
canonical settings are $d=4$, $d_{hi}=24$, eight heads, eight local attention
layers, $K=40$, voltage limits, direct updates, geometric-safe Armijo,
\texttt{mse\_weight}=5, graph-normalized log-cosh physics, signed-log residual
features, learning rate $10^{-5}$ and cosine annealing.  G1--G5 use the
37,022-row \texttt{GBnetwork} parquet.  The contemporaneous \texttt{case300}
G1--G5 runs use the 27,994-row comparison parquet; the separate 37,930-row
authoritative v2 G0 is not a paired baseline for those runs.

The external GraphKit/GridSFM rows below come from a different campaign: one
model per grid, seed 42, approximately 40 epochs, and the report's original
PPC-v2 data.  They are useful architectural references, but are not a strict
same-checkpoint comparison with the later G1--G5 campaign.

\section{GBnetwork results}

\begin{table}[H]
\centering\footnotesize
\begin{tabular}{lrrrr}
\toprule
Setting & $|V|$ RMSE & $\theta$ RMSE (deg) & mean case-max $|\Delta P|$ & mean case-max $|\Delta Q|$\\
\midrule
G0 signed-log reference & $\approx0.0302$ & $\approx0.889$ & not paired & not paired\\
G1 mean/max pre & 0.030015 & 1.1963 & 5.3509 & 108.51\\
\textbf{G2 attention pre} & \textbf{0.027269} & \textbf{0.9513} & \textbf{7.4611} & \textbf{68.119}\\
\textbf{G3 attention post} & \textbf{0.029390} & \textbf{1.0849} & \textbf{3.6316} & \textbf{46.803}\\
G4 gated pre & 0.030777 & 1.1578 & 7.1467 & 32.962\\
G5 gated post & 0.031904 & 1.3042 & 4.4431 & 62.240\\
G3 with MSE physics & 0.035169 & 1.4251 & 7.9892 & 205.45\\
GraphKit reference & 0.0354 & 4.94 & 181 & 55.8\\
\textbf{GridSFM reference} & \textbf{0.0319} & \textbf{2.14} & \textbf{1.10} & \textbf{4.81}\\
\bottomrule
\end{tabular}
\caption{GBnetwork voltage metrics and mean per-scenario maxima.  Every
residual entry in the last two columns uses
$N_s^{-1}\sum_s\max_i|\Delta P_{s,i}|$ (and analogously for $Q$); none is a
global test-set maximum.  External rows remain cross-campaign references.
``G3 with MSE physics'' is G3 with \texttt{physics\_loss\_form=mse} and nothing
else changed; it is a paired ablation of G3, not a sixth global-context
variant.}
\end{table}

\begin{table}[H]
\centering\small
\begin{tabular}{lrr}
\toprule
Setting & $R^2_{\sigma}$ on $|V|$ & $R^2_{\sigma}$ on $\theta$\\
\midrule
G1 mean/max pre & -0.611 & 0.971\\
\textbf{G2 attention pre} & \textbf{-0.329} & \textbf{0.982}\\
G3 attention post & -0.544 & 0.976\\
G4 gated pre & -0.694 & 0.973\\
G5 gated post & -0.820 & 0.965\\
G3 with MSE physics & -1.211 & 0.959\\
GraphKit reference & -1.241 & 0.502\\
GridSFM reference & -0.820 & 0.906\\
\bottomrule
\end{tabular}
\caption{GBnetwork coefficient of determination.  G1--G5 use the report's
reference spreads, $\sigma_{|V|}=0.02365$ pu and
$\sigma_\theta=7.0^\circ$, through
$R^2_{\sigma}=1-(\mathrm{RMSE}/\sigma)^2$.  GraphKit and GridSFM are the values printed
in the four-model report.}
\end{table}

\begin{table}[H]
\centering\footnotesize
\setlength{\tabcolsep}{2.2pt}
\begin{tabular}{lrrrrrrrrrr}
\toprule
Setting & mn P & mn Q & rms P & rms Q & p95 P & p95 Q & md P & md Q & $\le$1e-2 P & $\le$1e-3 P\\
\midrule
G1 & 0.1173 & 0.2953 & 0.3024 & 2.6399 & 0.4355 & 0.7875 & 0.0455 & 0.1333 & 19.94\% & 3.66\%\\
G2 & 0.1418 & 0.3056 & 0.3910 & 1.8777 & 0.5961 & 0.9023 & 0.0440 & 0.1336 & 20.87\% & 3.90\%\\
\textbf{G3} & \textbf{0.1094} & \textbf{0.2715} & \textbf{0.2477} & \textbf{1.3677} & \textbf{0.4332} & \textbf{0.7913} & \textbf{0.0411} & \textbf{0.1256} & \textbf{22.11\%} & \textbf{4.29\%}\\
G4 & 0.1225 & 0.2897 & 0.4447 & 1.0635 & 0.4388 & 0.9575 & 0.0478 & 0.1272 & 20.09\% & 3.67\%\\
G5 & 0.1114 & 0.2865 & 0.3246 & 1.5413 & 0.4069 & 0.8830 & 0.0468 & 0.1315 & 19.78\% & 3.66\%\\
G3-mse & 0.1252 & 0.3528 & 0.5837 & 4.8016 & 0.4754 & 0.8244 & 0.0452 & 0.1278 & 20.34\% & 3.85\%\\
\midrule
RANGE $M{=}1$ & 0.1095 & 0.2859 & 0.2695 & 2.7854 & 0.4104 & 0.7655 & 0.0465 & 0.1291 & 18.77\% & 3.43\%\\
RANGE $M{=}4$ & 0.2771 & 0.5772 & 1.0370 & 5.0045 & 1.1322 & 1.7964 & 0.0896 & 0.2015 & 12.99\% & 2.74\%\\
RANGE $M{=}8$ & 0.1696 & 0.6017 & 0.5973 & 4.8743 & 0.6326 & 2.0083 & 0.0649 & 0.1918 & 16.46\% & 4.23\%\\
\midrule
GraphKit & 2.05 & 0.771 & --- & --- & 6.85 & 2.74 & 0.458 & 0.191 & --- & 0.00\%\\
\textbf{GridSFM} & \textbf{0.0828} & \textbf{0.266} & --- & --- & \textbf{0.304} & \textbf{0.913} & \textbf{0.0432} & \textbf{0.156} & --- & \textbf{0.03\%}\\
\bottomrule
\end{tabular}
\caption{GBnetwork residual distributions, exact-complex test summary fields,
in pu on a 100-MVA base.  \textbf{Correction to the previous version of this
table:} the columns previously headed ``mean P'' and ``mean Q'' carried the
driver's \texttt{RMSE $\Delta$P}/\texttt{$\Delta$Q} fields, not the pooled
means, and the columns headed ``tol'' were the $\le10^{-3}$ pu fractions and
not the $10^{-6}$ pu tolerance.  Both are now separated and relabelled; no
underlying run changed.  The $10^{-6}$ pu fraction is $\le0.02\%$ for every
PIGNN row here and is omitted.  The external GraphKit/GridSFM rows are the
four-model report's printed mean/median/p95 fields; that report did not print
RMSE or the $10^{-2}$ pu fraction, so those cells are blank, and its ``tol''
column is used in the $\le10^{-3}$ position.}
\end{table}

The PIGNN G2 voltage RMSE is about 9.7\% below the historical G0 value.  G3
reduces the PIGNN mean case-max residuals relative to G2, especially reactive
power, but does not simultaneously minimize voltage error.

\section{case300 results}

\begin{table}[H]
\centering\footnotesize
\begin{tabular}{lrrrr}
\toprule
Setting & $|V|$ RMSE & $\theta$ RMSE (deg) & mean case-max $|\Delta P|$ & mean case-max $|\Delta Q|$\\
\midrule
G1 mean/max pre & 0.032798 & 4.9682 & 10.914 & 10.885\\
\textbf{G2 attention pre} & \textbf{0.027823} & \textbf{4.5222} & \textbf{5.463} & \textbf{7.260}\\
G3 attention post & 0.026900 & 4.5463 & 6.369 & 9.050\\
G4 gated pre & 0.032155 & 5.4597 & 10.160 & 11.566\\
G5 gated post & 0.034005 & 5.3795 & 8.932 & 11.859\\
\textbf{GraphKit reference} & \textbf{0.0237} & \textbf{4.76} & \textbf{10.4} & \textbf{6.82}\\
\textbf{GridSFM reference} & \textbf{0.0247} & \textbf{27.7} & \textbf{1.22} & \textbf{1.23}\\
\bottomrule
\end{tabular}
\caption{case300 comparison.  The G1--G5 rows use the 27,994-row campaign;
the external rows use the four-model report's case300 campaign.}
\end{table}

Global context does not provide a uniform case300 improvement.  G3 obtains
the lowest magnitude RMSE, while G2 obtains the lowest reported mean case-max
P/Q residuals.  The angle errors remain large for all five variants.  On GBnetwork the
per-bus/per-edge split in Section~\ref{sec:metrology} shows roughly a third to
two thirds of such an error to be gauge drift rather than flow error; the same
split has not been measured on case300, so the case300 angle column should be
read as an upper bound on the flow-relevant error rather than as its size.

\section{RANGE final test results}

The three RANGE training jobs were stopped by the four-hour Slurm wall clock
(state \texttt{TIMEOUT}) after 38, 33 and 33 of the planned 40 epochs for
$M=1,4,8$; their best checkpoints were written at epochs 35, 27 and 33
respectively.  Because the in-job test block never ran, the saved checkpoints
were re-evaluated in a separate test-only array (job 797689, three tasks, all
\texttt{COMPLETED}) on the same 12{,}344-scenario test split, seed 42, with the
identical flags.  The values below are therefore final test values, not
validation values, but they belong to slightly under-trained models.

\begin{table}[H]
\centering\footnotesize
\setlength{\tabcolsep}{3.5pt}
\begin{tabular}{lrrrrrrr}
\toprule
Masters & best ep. & $|V|$ RMSE & $\theta$ RMSE & $R^2_{\sigma,|V|}$ & $R^2_{\sigma,\theta}$ & mean c-max $|\Delta P|$ & mean c-max $|\Delta Q|$\\
\midrule
\textbf{$M=1$} & \textbf{35/38} & \textbf{0.03253} & \textbf{1.016$^\circ$} & \textbf{-0.891} & \textbf{+0.979} & \textbf{4.007} & \textbf{100.56}\\
$M=4$ & 27/33 & 0.04455 & 2.966$^\circ$ & -2.548 & +0.821 & 31.426 & 211.25\\
$M=8$ & 33/33 & 0.04683 & 1.604$^\circ$ & -2.921 & +0.947 & 17.843 & 205.35\\
\bottomrule
\end{tabular}
\caption{RANGE final test metrics.  ``best ep.'' is the epoch of the loaded
checkpoint over the epochs actually reached before the wall clock.  The
$R^2_{\sigma}$ values are recomputed here with the same reference spreads used for G1--G5,
$\sigma_{|V|}=0.02365$ pu and $\sigma_\theta=7.0^\circ$, so that this table can
be read against Table~2.  \textbf{These are not the driver's printed $R^2$
fields}, which for these three runs read $+0.610/+0.269/+0.193$ on $|V|$.  The
two differ only in the baseline predictor in the denominator, and the gap is
large.  \texttt{train\_valid\_test.py} pools every bus--scenario value and
subtracts a single global mean, so its baseline is ``predict one number
everywhere''; measured on this parquet that spread is
$\sigma_{|V|}^{\text{global}}=5.23\times10^{-2}$ pu, because it also contains
the \emph{static} bus-to-bus differences in nominal voltage level.  The spread
used throughout this report subtracts each bus's own mean first, so its
baseline is ``predict each bus's usual voltage'', giving
$\sigma_{|V|}^{\text{per-bus}}=2.38\times10^{-2}$ pu.  Only the latter measures
what the surrogate must actually learn --- how voltages move as the loading
changes --- since per-bus nominal levels are already available from
\texttt{vn\_kv} and the admittance matrix.  The GraphKit and GridSFM drivers
emit no $R^2$ at all, so every external $R^2$ in this report is likewise a
per-bus-baseline hand computation of $R^2_{\sigma}$; mixing the two conventions
in one column would make a model with a \emph{worse} $|V|$ RMSE appear to be
the only one with a positive value.}
\end{table}

\begin{table}[H]
\centering\small
\begin{tabular}{lrrrrr}
\toprule
Masters & agg./bcast.\ entropy & eff.\ masters & attn.\ cosine & emb.\ cosine & self-loop mass\\
\midrule
$M=1$ & 1.000 / 1.000 & 1.000 & +0.000 & +0.000 & 50.0\%\\
$M=4$ & 1.000 / 1.000 & 4.000 & +1.000 & +1.000 & 20.0\%\\
$M=8$ & 1.000 / 1.000 & 8.000 & +1.000 & +1.000 & 11.2\%\\
\bottomrule
\end{tabular}
\caption{RANGE master diagnostics on the test split.  The effective master
count equals the nominal count only because attention is uniform; for $M=4$
and $M=8$ the master-attention and master-embedding cosine similarities are
both $+1.000$, so the masters are duplicates rather than specialists.  Token
utilization is exactly uniform ($1/M$ each) in both cases.}
\end{table}

The final test values confirm the interim reading and sharpen it.  $M=1$ is the
best RANGE setting on every metric except pooled RMSE $\Delta Q$, and its
voltage RMSE (0.0325) and angle RMSE ($1.02^\circ$) place it between G1 (0.0300,
$1.20^\circ$) and G5 (0.0319, $1.30^\circ$) --- that is, \emph{inside} the
G1--G5 band rather than beyond it.  Its pooled mean and median residuals are
likewise indistinguishable from G3's (mean $|\Delta P|$ 0.1095 vs.\ 0.1094).
Adding masters is strictly harmful: $M=4$ roughly doubles every pooled residual
statistic and triples the angle error, and $M=8$ recovers only part of that.

The collapse diagnostics explain why.  With $M>1$ the masters converge to a
single shared representation, so the extra tokens add parameters and dilute the
per-token attention mass (self-loop mass falls from 50\% to 11\%) without
adding any distinct global channel.  A single master avoids the collapse
trivially, but then the sidecar reduces to the same one-token global summary
already tested as G2/G3 --- and it performs like it.  On this evidence,
persistent master tokens are not a productive direction for GBnetwork, and the
$M=1$ result should not be read as a positive RANGE finding.

One caveat: all three runs stopped 2--7 epochs short and $M=1$ was still
improving at the wall clock (validation loss $1.60\times10^{-2}$ at epoch 35
against $1.69\times10^{-2}$ at 38, without a clear plateau).  A rerun with a
longer wall clock could move these numbers, but not by the factor that would be
needed to separate $M=1$ from G2/G3.

\section{Where the residual mass actually is}
\label{sec:massdiag}

Every table above reduces the residual to a scalar, which cannot distinguish a
model that is uniformly slightly worse from one that is equal in the bulk and
catastrophic on a handful of buses.  A dedicated diagnostic (alex array
4123954/4124031, six tasks, all \texttt{COMPLETED}) re-evaluated the G1--G5 and
GridSFM v2 checkpoints on the identical 12{,}344-scenario test split, identical
P/Q masks and a final complex128 recomputation, and broke the residuals down by
bus type, nominal voltage, $|Y_{ii}|$, degree and transformer adjacency.

The rescore reproduces the published voltage metrics to within $0.8\%$ on all
six models ($|V|$ within $0.12\%$, worst angle deviation $-0.77\%$ on G3), so
the breakdowns below describe the same checkpoints as Tables 1--3 and not a
different evaluation.

\begin{table}[H]
\centering\footnotesize
\setlength{\tabcolsep}{3pt}
\begin{tabular}{lrrrrrrrr}
\toprule
& \multicolumn{4}{c}{$|\Delta P|$} & \multicolumn{4}{c}{$|\Delta Q|$}\\
\cmidrule(lr){2-5}\cmidrule(lr){6-9}
Setting & median & RMSE & glob.\ max & \%$>$1 pu & median & RMSE & glob.\ max & \%$>$10 pu\\
\midrule
G1 & 0.0455 & 0.302 & 36.1 & 1.03\% & 0.1333 & 2.642 & 172.4 & 0.057\%\\
G2 & 0.0442 & 0.398 & 37.4 & 2.11\% & 0.1338 & 1.878 & 157.2 & 0.074\%\\
G3 & \textbf{0.0410} & 0.244 & 18.6 & 1.00\% & \textbf{0.1255} & 1.354 & 141.4 & 0.061\%\\
G4 & 0.0478 & 0.465 & 304.2 & 1.19\% & 0.1273 & 1.069 & 230.2 & 0.062\%\\
G5 & 0.0468 & 0.324 & 58.8 & 0.79\% & 0.1315 & 1.539 & 141.9 & 0.061\%\\
GridSFM & 0.0432 & \textbf{0.141} & \textbf{2.4} & \textbf{0.05\%} & 0.1557 & \textbf{0.458} & \textbf{9.9} & \textbf{0.000\%}\\
\bottomrule
\end{tabular}
\caption{Bulk versus tail on GBnetwork, pu on a 100-MVA base.  ``glob.\ max''
is the true $\max_{s,i}|\Delta\cdot_{s,i}|$ over all bus--scenario entries, a
quantity none of the earlier tables reported.  \textbf{PIGNN is not broadly
worse than GridSFM}: G3's median residual is lower on both channels
(0.0410 vs.\ 0.0432 for P, 0.1255 vs.\ 0.1557 for Q) and its pooled mean
$|\Delta Q|$ is within 2\% of GridSFM's (0.271 vs.\ 0.266, Table~3).  The whole
gap is in the extreme tail: RMSE differs by $1.7$--$3.3\times$ and the global
maximum by one to two orders of magnitude.}
\end{table}

\subsection*{The tail is four buses}

Splitting by nominal voltage localizes it almost completely.  The 22-kV bin
contains $49{,}376$ entries $=4$ buses $\times\,12{,}344$ scenarios --- four
buses out of 2224, $0.18\%$ of the grid.

\begin{table}[H]
\centering\footnotesize
\begin{tabular}{lrrrr}
\toprule
Setting & mean $|\Delta Q|$ at 22 kV & \%$>$10 pu there & glob.\ max there & share of total $Q$ mass\\
\midrule
G1 & 27.71 & 24.8\% & 172.4 & \textbf{20.3\%}\\
G2 & 17.23 & 23.4\% & 157.2 & 12.2\%\\
G3 & 11.80 & 22.3\% & 141.4 & 9.4\%\\
G4 & 8.55 & 25.0\% & 230.2 & 6.4\%\\
G5 & 16.02 & 24.8\% & 141.9 & 12.1\%\\
\textbf{GridSFM} & \textbf{0.217} & \textbf{0.0\%} & \textbf{3.5} & \textbf{0.2\%}\\
\bottomrule
\end{tabular}
\caption{The four 22-kV buses.  ``Share of total $Q$ mass'' is
$\bar q_{22}n_{22}/(\bar q_{\text{pool}}n_{\text{pool}})$, the fraction of the
grid's entire pooled $|\Delta Q|$ contributed by these four buses.  For G1 they
carry a fifth of it.  The same buses cost GridSFM $0.2\%$.  The $P$ channel
shows the identical pattern more weakly ($1.6$--$5.0\%$ of $P$ mass for
G1--G5, $0.2\%$ for GridSFM).}
\end{table}

A paired G3 run with the physics loss changed from log-cosh to MSE
(\texttt{physics\_loss\_form=mse}, every other flag identical) is queued as
helma job 797773 to test whether a quadratic penalty suppresses this tail; a
first attempt on an alex A40 died in epoch~1 with a CUDA out-of-memory at the
attention logits, which is a GPU-capacity failure and not a property of the
objective.  Its result is not included here.

\subsection*{It tracks $|Y_{ii}|$, not the grid as a whole}

\begin{table}[H]
\centering\footnotesize
\begin{tabular}{lrrrrr}
\toprule
$|Y_{ii}|$ bin & entries & G3 mean $|\Delta Q|$ & G3 \%$>$10 pu & GridSFM mean & GridSFM \%$>$10\\
\midrule
$10^{0}$--$10^{1}$ & 4{,}073{,}520 & 0.153 & 0.000\% & 0.142 & 0.000\%\\
$10^{1}$--$10^{2}$ & 7{,}319{,}992 & 0.243 & 0.000\% & 0.205 & 0.000\%\\
$10^{2}$--$10^{3}$ & 7{,}221{,}240 & \textbf{0.244} & 0.000\% & 0.351 & 0.000\%\\
$10^{3}$--$10^{4}$ & 2{,}703{,}336 & \textbf{0.280} & 0.000\% & 0.322 & 0.000\%\\
$10^{4}$--$10^{5}$ & 1{,}234{,}400 & 0.738 & 0.892\% & 0.355 & 0.000\%\\
$10^{5}$--$10^{6}$ & 148{,}128 & 2.267 & 1.971\% & 0.953 & 0.000\%\\
\bottomrule
\end{tabular}
\caption{$|\Delta Q|$ against the magnitude of the diagonal admittance.  Below
$|Y_{ii}|=10^{4}$ PIGNN G3 matches or beats GridSFM and never exceeds 10 pu;
above it, G3 degrades by $3$--$9\times$ while GridSFM does not.  The degree
breakdown agrees: the single degree-15 bus has G3 mean 6.87 pu with $8.4\%$ of
entries above 10 pu, against GridSFM's 0.474 pu and $0.0\%$.}
\end{table}

Two conclusions follow, and both revise the reading in the next section.

First, \textbf{the PIGNN deficit is a stiffness-localized tail, not a
representation-capacity deficit}.  The failure sits exactly on the
worst-conditioned nodes, is confined to well under $1\%$ of entries, and leaves
the bulk distribution competitive or better.  This is the failure mode a
diagonal equilibration of the residual features is meant to address, and it is
the clearest evidence so far that the conditioning hypothesis --- rather than
receptive field or parameter count --- explains the large-grid gap.

Second, \textbf{GridSFM's advantage looks like a bound, not better typical
accuracy}.  Its \%$>$10 pu is exactly $0.000$ in \emph{every} bin of every
breakdown --- bus type, voltage level, $|Y_{ii}|$, degree and transformer
adjacency alike --- while its median $|\Delta Q|$ is the worst of the six
models.  A model that is uniformly capped and slightly loose in the middle is
what a clipped magnitude head and a bounded angle residual around a DC prior
would produce.  The corresponding PIGNN mechanism, the geometric-safe Armijo
line search, bounds the \emph{step} rather than the \emph{output}, which does
not prevent a badly-conditioned bus from accumulating error over $K=40$
iterations.

\section{Notation: three different $R^2$, and the calibration}
\label{sec:r2notation}

Three distinct coefficients appear in the literature and in this project's
logs, and earlier revisions of this report used the word ``$R^2$'' for more
than one of them.  Every table below names which one it is.  Write $v$ for the
reference $|V|$ at a bus in a scenario, $\hat v$ for the prediction,
$\bar v_i$ for bus $i$'s mean over scenarios, and fit
$\hat v = a\,v + b + \varepsilon$ by least squares.

\begin{table}[H]
\centering\footnotesize
\setlength{\tabcolsep}{4pt}
\begin{tabular}{llll}
\toprule
Symbol & Definition & Null model it scores $0$ against & Calibration\\
\midrule
$R^2_{\mathrm{pool}}$ & Pearson $R^2$ of $\hat v$ on $v$, pooled over
  & one number for every bus & changes\\
  & all buses and scenarios & and scenario & \\[2pt]
$R^2_{\mathrm{bus}}$ & the same after subtracting $\bar v_i$
  & each bus at its own mean & \textbf{invariant}\\
  & from $\hat v$ and $v$ alike & & \\[2pt]
$R^2_{\sigma}$ & $1-\sum(\hat v-v)^2/\sum(v-\bar v_i)^2$
  & each bus at its own mean & changes\\
\bottomrule
\end{tabular}
\end{table}

$R^2_{\mathrm{pool}}$ is the diagnostic GridSFM's section 6.2 defines, and it
is the right one for a multi-grid foundation model, which must reproduce a
static voltage profile it has never seen.  Every campaign here trains one model
per grid on that grid, where that profile is memorisable from the training
split; a predictor that emits each bus's training mean and ignores the scenario
entirely scores $R^2_{\mathrm{pool}} = 0.798$ on GBnetwork.  So
$R^2_{\mathrm{bus}}$ is the tracking question in this setting, and
$R^2_{\mathrm{pool}}$ is kept only for comparability with the published
numbers.

$R^2_{\sigma}$ is what earlier revisions of this report printed.  It is not a
tracking measure: it charges the model for bias and for amplitude error as
well.  Writing $\sigma^2=\tfrac{1}{N}\sum(v-\bar v_i)^2$ and
$\sigma_\varepsilon^2$ for the variance of $\varepsilon$, the mean square error
splits exactly into three terms,
\[
  \mathrm{RMSE}^2 \;=\; \underbrace{b^2}_{\text{bias}}
  \;+\;\underbrace{(a-1)^2\sigma^2}_{\text{amplitude}}
  \;+\;\underbrace{\sigma_\varepsilon^2}_{\text{noise}},
  \qquad\text{so}\qquad
  R^2_{\sigma} \;=\; 1-\frac{b^2}{\sigma^2}-(a-1)^2-\frac{\sigma_\varepsilon^2}{\sigma^2}.
\]
$R^2_{\mathrm{bus}}$ normalises the noise term by the variance of the
\emph{prediction} rather than of the reference, so it is not a simple
difference away from $R^2_{\sigma}$; the two answer different questions and
should not be converted into one another.

Three relations are used below, and each was checked against the measurements
rather than assumed.  (i) Calibration removes exactly the $b^2/\sigma^2$ term,
so $R^2_{\sigma,\mathrm{cal}} = 1-(a-1)^2-\sigma_\varepsilon^2/\sigma^2$; over
the 26 stabilization runs this held to $1.4\times10^{-2}$.  (ii)
$R^2_{\mathrm{bus}}$ and the slope are unchanged by calibration, since a
per-bus offset cannot alter a regression on per-bus deviations; measured
deviation $<10^{-6}$ on all 26.  (iii) $R^2_{\sigma,\mathrm{cal}} \le
R^2_{\mathrm{bus}}$, with the gap being the amplitude error; observed on all
26, from $0.0006$ to $1.22$.

\paragraph{Post-hoc per-bus calibration.}  The $b^2/\sigma^2$ term is removable
without touching the model.  For each bus, take the mean prediction over the
\emph{training} split, align it to the training reference mean, and subtract
that fixed offset from every unseen test scenario:
\[
  \hat v^{\,\mathrm{cal}}_{i,s} \;=\; \hat v_{i,s}
  \;-\;\bigl(\overline{\hat v}^{\,\mathrm{train}}_{i}-\bar v^{\,\mathrm{train}}_{i}\bigr),
\]
and the circular equivalent for $\theta$.  No test data enters the fit, so this
is a correction a deployed model could carry.  Because it is fitted on one
split and applied to another it does not transfer perfectly: across the 26
stabilization runs the realised improvement in $R^2_{\sigma}$ differs from the
fitted $b^2/\sigma^2$ by a median of $1.2\%$ and at most $8.4\%$.  Tables that
report it use the form \emph{raw $\to$ calibrated}.  $R^2_{\mathrm{bus}}$ and the slope are
mathematically unchanged by it --- a per-bus offset cannot alter a regression
on per-bus deviations --- which is used below as a consistency check on the
implementation.

\section{What the metrics were measuring}
\label{sec:metrology}

Every ranking above rests on $|V|$ RMSE, $\theta$ RMSE and $R^2$.  Three
follow-up runs on the same checkpoints, split and masks show that each of the
three was answering a different question from the one it was being read as.

\paragraph{A predictor that ignores the scenario.}  The natural null model here
is not $V=1.0$ everywhere but ``each bus sits at its own usual voltage'': take
each bus's mean $|V|$ and circular-mean $\theta$ over the \emph{training} third
and emit them regardless of loading.  Scored on the same test third (alex job
4124582) it reproduces the reference spreads exactly --- $|V|$ RMSE
$2.3412\times10^{-2}$ against a measured
$\sigma_{|V|}^{\text{per-bus}}=2.3412\times10^{-2}$ pu, $\theta$ RMSE
$6.846^\circ$ against $6.846^\circ$ --- which is the definitional check that it
scores $R^2=0$ on both channels.

\begin{table}[H]
\centering\footnotesize
\setlength{\tabcolsep}{4pt}
\begin{tabular}{lrrrrrr}
\toprule
Setting & $|V|$ RMSE & $\theta$ RMSE & med $|\Delta P|$ & med $|\Delta Q|$ & RMSE $\Delta Q$ & max $|\Delta Q|$\\
\midrule
\textbf{per-bus constant} & \textbf{0.02341} & 6.846$^\circ$ & \textbf{0.0033} & \textbf{0.0015} & \textbf{0.088} & \textbf{3.4}\\
\midrule
G1 & 0.03002 & 1.197$^\circ$ & 0.0455 & 0.1332 & 2.642 & 172.3\\
G2 & 0.02727 & \textbf{0.946}$^\circ$ & 0.0441 & 0.1337 & 1.874 & 139.1\\
G3 & 0.02940 & 1.073$^\circ$ & 0.0410 & 0.1255 & 1.355 & 141.4\\
G4 & 0.03078 & 1.184$^\circ$ & 0.0478 & 0.1272 & 1.076 & 249.2\\
G5 & 0.03191 & 1.302$^\circ$ & 0.0468 & 0.1314 & 1.539 & 141.9\\
GridSFM & 0.03186 & 2.144$^\circ$ & 0.0432 & 0.1557 & 0.458 & 9.9\\
G3-mse & 0.03517 & 1.425$^\circ$ & 0.0452 & 0.1278 & 4.802 & 215.6\\
\bottomrule
\end{tabular}
\caption{Against a model that has learned nothing about the scenario.  It wins
on $|V|$ RMSE, on both median residuals, on pooled RMSE $\Delta Q$ and on the
worst bus in the corpus.  The only column any trained model wins is the bus
angle.  Residuals in pu on a 100-MVA base.}
\end{table}

\paragraph{On median residuals against the constant predictor.} On GBnetwork
$76.8\%$ of the $Q$-scored buses have $Q^{\mathrm{set}}$ exactly zero, and
$64.6\%$ of $P$-scored buses have $P^{\mathrm{set}}$ exactly zero; the median
setpoint is $0.0000$~pu on both channels. A median residual therefore lands on
pure transit nodes, where a predictor that emits a valid AC solution of the
mean scenario is exact for free --- no injection, no mismatch, in any scenario
--- while a learned model perturbs the voltage everywhere and gets no such
gift. Median comparisons \emph{between models} are unaffected, since the
distortion is common to them; median comparisons \emph{against the constant
predictor} are not a measure of skill. The mean and the maximum are dominated
by the buses that carry injection and are the honest comparison there.

\paragraph{The angle column does not mean what it appears to.}  Branch flows
depend on $\theta_i-\theta_j$, never on the absolute bus angle, so a common
per-scenario offset inflates per-bus $\theta$ RMSE without moving a single flow.
Splitting the error accordingly (deg):

\begin{table}[H]
\centering\footnotesize
\begin{tabular}{lrrrr}
\toprule
Setting & per-bus & of which gauge & de-gauged & \textbf{per-edge $\Delta\theta_{ij}$}\\
\midrule
per-bus constant & 6.846 & 4.029 & 5.532 & \textbf{1.000}\\
G2 & 0.946 & 0.330 & 0.887 & \textbf{0.376}\\
G3 & 1.073 & 0.483 & 0.958 & 0.403\\
G5 & 1.302 & 0.837 & 0.997 & 0.425\\
GridSFM & 2.144 & 1.575 & 1.454 & 0.447\\
\bottomrule
\end{tabular}
\caption{Angle error RMSE.  ``gauge'' is the circular mean of the bus-angle
error within each scenario.  On per-bus RMSE G2 beats the constant predictor by
$7.2\times$ and GridSFM by $2.3\times$; on the quantity the flows see those
margins fall to $2.7\times$ and $1.2\times$.  More than two thirds of GridSFM's
angle error is gauge, which is the criticism its own paper levels at
gridfm-graphkit.  The angle ranking in Table~1 and Table~5 is therefore
substantially a ranking of gauge drift.}
\end{table}

\paragraph{Two different $R^2$, and only one of them is about tracking.}
GridSFM's section 6.2 diagnostic is a pooled regression of prediction against
reference, reported with a slope; its $R^2$ is a Pearson $R^2$ and asks whether
the model orders scenarios correctly.  The $R^2$ in Tables 2 and 5 of this
report is $1-(\mathrm{RMSE}/\sigma)^2$, which also charges the model for bias
and for amplitude error.  They coincide only when the slope is 1 and the bias
is zero.  Furthermore, pooling over buses is right for GridSFM --- a multi-grid
foundation model must predict a grid's static profile because it has not seen
it --- but wrong here, where one model is trained per grid and that profile is
memorisable from the training split.

\begin{table}[H]
\centering\footnotesize
\begin{tabular}{lrrrrr}
\toprule
& \multicolumn{2}{c}{pooled} & \multicolumn{2}{c}{per-bus} & \\
\cmidrule(lr){2-3}\cmidrule(lr){4-5}
Setting & slope & $R^2_{\mathrm{pool}}$ & slope & $R^2_{\mathrm{bus}}$ & $R^2_{\sigma}$\\
\midrule
per-bus constant & 0.798 & 0.798 & \textbf{0.000} & \textbf{0.000} & $-0.000$\\
\midrule
G1 & 0.587 & 0.681 & 0.487 & 0.621 & $-0.644$\\
G2 & 0.640 & 0.740 & 0.648 & 0.709 & $-0.357$\\
G3 & 0.578 & 0.705 & 0.491 & 0.686 & $-0.577$\\
G4 & 0.554 & 0.676 & 0.510 & 0.652 & $-0.728$\\
G5 & 0.539 & 0.643 & 0.476 & 0.592 & $-0.857$\\
GridSFM & 0.575 & 0.631 & \textbf{0.688} & \textbf{0.752} & $-0.852$\\
G3-mse & 0.629 & 0.611 & 0.316 & 0.513 & $-1.258$\\
\bottomrule
\end{tabular}
\caption{$R^2$ on $|V|$ under the three readings of
Section~\ref{sec:r2notation}.  The first row is the
control: a predictor that reproduces each bus's mean and nothing else scores
$0.798$ pooled --- four fifths of the pooled variance on this grid is the
static profile, not the scenario.  \textbf{Every model retains
$R^2=0.59$--$0.75$ after that profile is removed}, so none of them is collapsed
in the sense the report previously claimed; the negative SSE column is measuring
something else.}
\end{table}

\paragraph{What the negative column is measuring: a per-bus offset.}  Writing
the prediction as $\hat v = a\,v + b + \varepsilon$ with $\varepsilon$
uncorrelated with $v$,
\[
  \mathrm{RMSE}^2 \;=\; \underbrace{b^2}_{\text{bias}}
  \;+\;\underbrace{(a-1)^2\sigma^2}_{\text{amplitude}}
  \;+\;\underbrace{\sigma_\varepsilon^2}_{\text{noise}},
  \qquad \sigma_\varepsilon^2=(1-R^2)\,\sigma_{\hat v}^2 .
\]

\begin{table}[H]
\centering\footnotesize
\begin{tabular}{lrrrr}
\toprule
Setting & $|V|$ RMSE & bias & amplitude & noise\\
\midrule
per-bus constant & 0.02341 & 0.00006 & 0.02341 & 0.00000\\
\midrule
G2 & 0.02727 & \textbf{0.02410} & 0.00825 & 0.00972\\
G3 & 0.02940 & \textbf{0.02572} & 0.01191 & 0.00779\\
G5 & 0.03191 & \textbf{0.02796} & 0.01227 & 0.00924\\
GridSFM & 0.03186 & \textbf{0.02960} & 0.00730 & 0.00926\\
G3-mse & 0.03517 & \textbf{0.02961} & 0.01602 & 0.00981\\
\bottomrule
\end{tabular}
\caption{Decomposition of the $|V|$ error, pu.  The control row validates it:
for a constant predictor the bias vanishes, the amplitude term is the full
$\sigma$ and there is no noise, so $\mathrm{RMSE}=\sigma$ exactly.  For every
trained model the \emph{bias} term alone exceeds the entire scenario spread
$\sigma=0.02341$: each bus's predicted mean voltage is systematically wrong by
more than the whole quantity being predicted.}
\end{table}

\paragraph{Removing the offset, and what it does not fix.}  This is a claim
about a correction learnable without the test set, so it was fitted where the
model's own training data lives --- a per-bus offset measured on the training
third, subtracted at test time, nothing retrained (alex job 4124803).

\begin{table}[H]
\centering\footnotesize
\setlength{\tabcolsep}{2.2pt}
\begin{tabular}{lrrrrrrr}
\toprule
& \multicolumn{3}{c}{$|V|$} & \multicolumn{2}{c}{residual medians} & \multicolumn{2}{c}{}\\
\cmidrule(lr){2-4}\cmidrule(lr){5-6}\cmidrule(lr){7-8}
Setting & RMSE & $R^2_{\sigma}$ & (pred.) & $|\Delta P|$ & $|\Delta Q|$ & $\Delta\theta_{ij}$ & max $|\Delta Q|$\\
\midrule
G1 & 0.0300$\to$\textbf{0.0150} & $-0.644\to$\textbf{$+0.592$} & $+0.592$ & 0.046$\to$0.027 & 0.133$\to$0.044 & 0.41$\to$0.29 & 172$\to$138\\
G2 & 0.0273$\to$\textbf{0.0128} & $-0.357\to$\textbf{$+0.703$} & $+0.704$ & 0.044$\to$0.030 & 0.134$\to$0.042 & 0.38$\to$0.27 & 139$\to$\textbf{197}\\
G3 & 0.0294$\to$\textbf{0.0142} & $-0.577\to$\textbf{$+0.630$} & $+0.630$ & 0.041$\to$0.024 & 0.126$\to$0.037 & 0.40$\to$0.28 & 141$\to$\textbf{160}\\
G4 & 0.0308$\to$\textbf{0.0144} & $-0.728\to$\textbf{$+0.621$} & $+0.621$ & 0.048$\to$0.026 & 0.127$\to$0.040 & 0.41$\to$0.29 & 249$\to$\textbf{357}\\
G5 & 0.0319$\to$\textbf{0.0154} & $-0.857\to$\textbf{$+0.569$} & $+0.569$ & 0.047$\to$0.025 & 0.131$\to$0.042 & 0.43$\to$0.30 & 142$\to$88\\
GridSFM & 0.0319$\to$\textbf{0.0118} & $-0.852\to$\textbf{$+0.746$} & $+0.747$ & 0.043$\to$0.014 & 0.156$\to$0.039 & 0.45$\to$0.23 & 9.9$\to$9.7\\
G3-mse & 0.0352$\to$\textbf{0.0176} & $-1.258\to$\textbf{$+0.438$} & --- & 0.045$\to$0.019 & 0.128$\to$0.031 & 0.40$\to$0.28 & \textbf{216$\to$56}\\
\bottomrule
\end{tabular}
\caption{Before and after subtracting a per-bus offset fitted on the training
split.  The ``predicted'' column is what the decomposition above says should
happen; measurement agrees to within $0.001$ on all six, which is the check
that the decomposition is the right one.  $|V|$ RMSE improves
$2.0$--$2.7\times$, every model passes the constant predictor for the first
time, median $|\Delta Q|$ improves $3.1$--$4.0\times$, and per-edge
$\Delta\theta_{ij}$ improves about $30\%$.  \textbf{The extreme tail does not
follow}: three of the six report a \emph{larger} worst-bus $|\Delta Q|$
afterwards.}
\end{table}

\paragraph{Log-cosh against MSE in the physics term.}  Section~\ref{sec:massdiag}
localised the PIGNN deficit to a handful of very large residuals, which makes a
quadratic penalty the obvious thing to try: MSE weights a residual of $10$ pu a
hundred times more than one of $1$ pu, where log-cosh weights it about ten
times.  G3 was therefore rerun with \texttt{physics\_loss\_form=mse} and every
other flag unchanged (helma job 797773, 40/40 epochs; a first attempt on an
alex A40 died in epoch 1 with a CUDA out-of-memory, which is a GPU-capacity
failure and not a property of the objective).

\textbf{It is worse on every axis, and worst exactly where it was meant to
help.}  Voltage RMSE rises from $0.02940$ to $0.03517$, the angle from
$1.073^\circ$ to $1.425^\circ$, $R^2_{\mathrm{bus}}$ falls from $0.686$ to
$0.513$ and the slope from $0.491$ to $0.316$ --- so the quadratic term costs
tracking as well as accuracy.  The worst bus goes the wrong way by an order of
magnitude: $\max|\Delta P|$ $18.6 \to 197.8$ pu, $\max|\Delta Q|$
$141.4 \to 215.6$ pu.  This is the failure mode the choice of log-cosh was
guarding against: with a squared penalty the few stiff buses of
Section~\ref{sec:massdiag} dominate the gradient and the fit is spent on them.

The one place it does better is under calibration, and revealingly so.  G3-mse
carries the largest \emph{amplitude} error of any setting here ($0.01602$ pu
against $0.00730$--$0.01227$ for the rest) and an ordinary bias, and once the
bias is removed its pooled $\Delta Q$ collapses --- RMSE $4.802 \to 0.221$, a
$22\times$ improvement, and $\max|\Delta Q|$ $215.6 \to 56.1$ --- whereas
calibrating G3 with log-cosh leaves its worst bus slightly \emph{worse}
($141 \to 160$).  Both end up in the same place: calibrated, G3-mse reaches
$R^2_{\sigma} = +0.438$ against G3's $+0.630$.  MSE is not a useful
substitute, but its damage is more of the kind a post-hoc offset can undo.

Three conclusions, and the first two overturn statements made earlier in this
report.

\textbf{The models are not failing to track the scenario.}  After the static
profile is removed they still explain 59--75\% of the variance of $|V|$.  What
they get wrong is each bus's mean level, by more than the spread they are
predicting, and that miscalibration is what every negative $R^2$ in this report
and in the four-model report is dominated by.  It is also free to remove: no
retraining, no architecture change, a $2$--$2.7\times$ reduction in $|V|$ RMSE
and $3$--$4\times$ in median $|\Delta Q|$.

\textbf{The angle advantage is real but roughly a third of its apparent size},
and for GridSFM it nearly vanishes: $2.144^\circ$ per bus against $0.447^\circ$
per edge.  Comparisons on per-bus $\theta$ RMSE across models with different
gauge behaviour are not meaningful.

\textbf{Neither correction touches the failure in
Section~\ref{sec:massdiag}.}  The worst-bus $|\Delta Q|$ is unchanged or worse
after debiasing, and the constant predictor still beats every debiased model on
median residual by an order of magnitude.  The stiffness-localized tail on the
four 22-kV buses is a separate defect from the calibration one, and only the
calibration defect has a free fix.

\section{Interpretation and bottleneck}

The GridSFM architecture offers plausible explanations for its large-grid
feasibility advantage.
GridSFM uses seven heterogeneous node types (bus, generator proxy, load,
shunt, line, transformer and cycle), promotes branches to nodes, includes a
fundamental cycle basis, and performs global linear attention within each node
type.  It also supplies DC/Laplacian features, effective-resistance landmarks,
branch loading features and explicit bus-type encodings.  Its angle head is a
bounded residual around a DC prior and its magnitude head is clipped to voltage
limits.  However, the ppNR-v2 run in the four-model report did \emph{not} use
the released native GridSFM multi-term objective.  Its dispatch used voltage
MSE plus a shared PPC masked log-cosh residual loss with physics weight
$10^{-2}$.  Native GridSFM loss components are therefore design references,
not an explanation directly tested by that run.

The current PIGNN G0 backbone instead performs repeated learned corrections
over a bus-centric one-hop graph.  Eight local layers provide a large formal
receptive field, but long-range information must still traverse local
nonlinearities.  The G2/G3 sidecars add one graph token; they do not reproduce
GridSFM's node-level same-type global attention.  PIGNN also composes update
errors over 40 steps, so a correction that improves voltage MSE can increase
the AC residual.  This is the principal bottleneck, rather than simply an
insufficient value of $K$.

The breakdown in Section~\ref{sec:massdiag} narrows this list considerably.
Since PIGNN already matches or beats GridSFM in the bulk and on every
$|Y_{ii}|$ bin below $10^{4}$, the ingredients that would explain a broad
representational advantage are not the ones being exercised.  What separates
the two models is confined to the worst-conditioned $0.2\%$ of buses, and
GridSFM's uniform absence of any entry above 10 pu points at its output bounds.
The two candidates worth testing are therefore \emph{bounded terminal
corrections} --- an output clip, as distinct from PIGNN's Armijo step bound ---
and \emph{diagonal equilibration of the residual features}, which targets the
$|Y_{ii}|$ dependence directly.  Same-type global communication and explicit
branch/cycle states remain plausible but are not supported by this evidence:
the global-context sidecars G1--G5 span a $2\times$ range in tail mass without
moving the bulk at all.  A
two-hop block is a useful intermediate test because it provides a sparse
mesoscopic shortcut, but it should be evaluated against a parameter-matched
one-hop control and then combined with RANGE-$M=1$ only if the standalone
two-hop result is positive.

\section{Metric and campaign caveats}

The historical field names \texttt{dPinf}/\texttt{dQinf} are potentially
misleading.  In both the four-model driver and the current PIGNN driver, the
per-scenario bus maximum is accumulated and divided by the scenario count:
\[
  \operatorname{case\text{-}max\ mean}(P)
  =\frac{1}{N_s}\sum_{s=1}^{N_s}\max_{i\in\mathcal M_P}|\Delta P_{s,i}|.
\]
The true global maximum $\max_{s,i}|\Delta P_{s,i}|$ was not reported by the
four-model tables; it is now measured for GBnetwork in
Section~\ref{sec:massdiag}, and it separates the models far more sharply than
the mean-of-maxima reducer does (G4 reaches 304 pu on $P$ where its mean
per-scenario maximum is only 7.3).  Pooled mean, median and quantiles instead reduce the full
set of valid bus--scenario entries; they must not be called case maxima.  A
strict head-to-head claim still requires the same split, masks and precision,
but not because the historical rows used a different maximum reducer.

The GraphKit/GridSFM campaign also predates the later PIGNN signed-log,
global-context and RANGE experiments.  Its results should therefore be used
as architectural evidence and motivation, not as a strict controlled
head-to-head estimate of the latest PIGNN gap.

\section{Recommended next step}

The immediate decision path is:
\begin{enumerate}
\item RANGE-$M=1,4,8$ is finished and test-rescored; treat it as closed and do
      not extend it to further master counts;
\item compare the two-hop H0--H5 campaign using both mean scenario maxima and
      the true global P/Q maxima, which Section~\ref{sec:massdiag} shows are the
      discriminating statistic;
\item retain G2 as the voltage-oriented global candidate and G3 as the
      feasibility-oriented candidate;
\item test the two mechanisms that Section~\ref{sec:massdiag} actually
      implicates --- a bounded output head, and diagonal
      ($|Y_{ii}|$-equilibrated) residual features --- on the four 22-kV buses
      before spending further effort on global-context topology;
\item if H2 or H5 is positive, run three seeds on the selected GBnetwork
      candidate before extending to additional grids.
\end{enumerate}

\end{document}
